\section{Introduction}
\label{sec:intro}

Modern agents write substantial programs. A single agent run now sequences many gen calls, branches on intermediate results, accumulates mutable state, and delegates subtasks to other agents. The volume of code an agent authors grows with every model generation. What does not grow is safety. A workflow of ten gen calls receives no more type discipline than a one-shot prompt, and authorization, when enforced at all, is applied at the outer boundary rather than at each call. The gap shows up on the two axes that matter. Gen outputs flow from one call to the next with no check that their shape matches what downstream code expects, and sensitive gens can be invoked from within a workflow with no evidence that the invocation was authorized. As agents are trusted with more consequential work such as filling forms, querying sensitive records, and generating other agents, safety must scale with the size of the workflow rather than lag behind it.

Our position is simple. Agents plan, and plans are programs. If plans are programs, then the mature body of programming-language techniques applies directly to them. Type systems discipline the data flowing between gen calls. Gradual typing~\cite{siek2006gradual} accommodates the values, types, workflows, and policies that only materialize at runtime when the LLM produces them. Static analysis catches structural errors before any LLM call is paid for, and authorization semantics let each gen invocation be checked against an explicit principal. The same type and gen-call information that constrains what the LLM may produce also shapes how it generates, so the infrastructure that catches errors is the infrastructure that prevents them. Safety that scales with the size of the workflow is exactly what these techniques deliver. Every new gen call, every new field access, and every new principal is another point at which the program can be checked.

Today developers express these plans using orchestration frameworks such as LangChain~\cite{langchain2022} or ad-hoc Python scripts. These approaches treat plans as configurations of library objects rather than as programs in their own right. Gen calls return untyped data, control flow is scattered across callback edges, structural errors surface only at runtime, and authorization is applied at the outer boundary rather than per invocation. Neighboring programming-language approaches such as LMQL~\cite{beurer2023prompting} and DSPy~\cite{khattab2023dspy} address parts of this problem but do not integrate a type system for gen outputs with per-invocation authorization. Yue et al.~\cite{yue2026statictemplatesdynamicruntime} identify the tension between expressivity and verifiability as a central open problem in this space.

\begin{figure}[H]
\begin{lstlisting}[language=hadl,showstringspaces=false]
@retries 3

workflow clinical_trial() -> String {
  let trial  = ask "Describe the clinical trial"
  let crf: Schema   = gen "Generate a CRF schema for `trial`"
  let policy: Policy = gen "Coordinators cannot see data in
    unblinded folder. Safety board has full access."
  enforce policy
  let num_visits = ask "How many patient visits to simulate?"
  let pids: Schema = gen "Return `num_visits` patient IDs"
  var cost: Number = 0
  for pid in pids {
    let visit: crf = gen "Generate a visit for `pid` in `trial`
      conforming to schema `crf`" as "coordinator"
    cost = js cost + visit.cost
    say "Visit `visit.patient_id`: cost `visit.cost`"
  }
  agent "Write a report for the trial visits" as "safety_board"
}
\end{lstlisting}
\caption{Clinical-trial workflow demonstrating gradual schema typing (\texttt{crf} typed binding),
  gradual policy typing (\texttt{enforce policy}), and per-invocation Cedar authorization (\texttt{as} clause).
  \texttt{@retries 3} activates feedback-driven self-healing: if the generated schema lacks a field
  accessed later (\texttt{visit.cost}, \texttt{visit.patient\_id}), the runtime retries the causal
  \texttt{gen} with the violated constraint as feedback before any visit data is collected.}
\label{fig:clinical-trial}
\end{figure}


Figure~\ref{fig:clinical-trial} illustrates gradual validation and per-invocation authorization in a single workflow. Two \texttt{gen} calls produce program structure, a type used in later bindings and a policy activated via \texttt{enforce} on the next line. Every subsequent \texttt{gen} and \texttt{agent} call is typed against the generated schema and checked against the generated policy under an explicit principal, including the closing \texttt{agent} call that returns an LLM-written summary. The same type discipline that governs ordinary field access such as \texttt{visit.patient\_id} governs the LLM output that produced the value.

We present HADL (Hierarchical Agent Dataflow Language), a strongly typed imperative domain-specific language for composing long-running agentic workflows. HADL provides four domain-specific primitives, \texttt{ask}, \texttt{say}, \texttt{gen}, and \texttt{agent}, as first-class typed expressions in an imperative language with mutable state, unbounded loops, structured iteration, and nested workflow composition. The type system supports three forms of gradual validation. The \texttt{Schema} type lets a gen produce a type that resolves to a concrete type once the LLM delivers one, after which subsequent uses are checked under the same type system. Workflow types let a gen produce a complete workflow that is then invoked via \texttt{eval} and checked at the call site. The \texttt{Policy} type lets a gen produce a Cedar~\cite{cutler2024cedar} policy that is then activated via \texttt{enforce}, with safety checks and append-only merging so that activation can only tighten access. The language and runtime are implemented in Rust, integrating with the GitHub Copilot SDK for LLM execution.

This paper makes two contributions.

\begin{enumerate}
	\item We design HADL, an imperative DSL in which \texttt{ask}, \texttt{say}, \texttt{gen}, and \texttt{agent} are first-class typed expressions with curried function types, providing natural imperative syntax for multi-turn agent orchestration. A type checker with source-span diagnostics catches structural errors before any LLM call is made. The language is implemented in Rust with a PEG parser, type checker, and async tree-walking interpreter.
	\item We introduce \emph{gradual validation}, a mechanism in which types, workflows, and policies are produced by LLMs and checked at runtime. Beyond catching errors, the declared types guide the LLM's generation itself, so every call produces output shaped to fit the surrounding program. \texttt{Schema} resolves to a concrete type once an LLM produces one, after which subsequent uses are checked under the same type system. Workflow types enable LLM-generated workflows executed via \texttt{eval}. \texttt{Policy} enables LLM-generated Cedar authorization policies activated via \texttt{enforce}, with a deny-all guard, scoped principal validation, and append-only merging. All three share provenance-driven self-healing that reruns the generator when validation fails.
\end{enumerate}

The remainder of this paper is organized as follows. Section~\ref{sec:overview} presents the language through motivating examples. Section~\ref{sec:design} presents the type system and gradual validation. Section~\ref{sec:semantics} gives the operational semantics and states soundness. Section~\ref{sec:implementation} describes the compilation pipeline and runtime. Section~\ref{sec:related} discusses related work, and Section~\ref{sec:conclusion} concludes.
